% sec-01-fda-forms.tex - FDA Forms 1571 and 3674 (final stage)
\section{FDA Forms 1571 and 3674}

The administrative gateway to this initial Investigational New Drug
(IND) application is the pair of executed Food and Drug Administration
(FDA) forms that open every Center for Drug Evaluation and Research (CDER)
submission: FDA Form 1571 (Investigational New Drug Application) and FDA
Form 3674 (Certification of Compliance, under 42 U.S.C. \S 282(j)). These
two forms are positioned, together with the Cover Letter, ahead of the
Table of Contents so that a reviewer encounters the sponsor agreement, the
serial number, the contents map, and the clinicaltrials.gov certification
before any scientific module. They are the regulatory envelope around the
LLM-Directed PDAC Robotic Daraxonrasib Phase 1 protocol (protocol
v1.0.0), a first-in-human, prospective, open-label, single-arm,
early-feasibility study that combines an IND under 21 CFR Part 312 with an
Investigational Device Exemption (IDE) framework under 21 CFR Part 812 for
an on-premises, LLM-directed, eight-arm robotic pancreaticoduodenectomy
(Whipple) device used together with oral daraxonrasib (RMC-6236), an oral
RAS(ON) multi-selective (pan-RAS) inhibitor \cite{daraxwhipple,
onpremwhippl}. This section summarizes the field-by-field content and the
governing logic of each form; the executed, signed PDF versions of FDA
Form 1571 and FDA Form 3674 are assembled with the physical submission and
are not reproduced inside this LaTeX source, in keeping with FDA's
expectation that the original signed forms accompany the package
\cite{phase1guidance, cfr312paiuni}.

The two forms discharge two distinct legal functions that together
constitute the administrative predicate for FDA's acceptance of the file.
Form 1571 is the instrument through which the sponsor-investigator,
ChemicalQDevice (Kevin Kawchak, Chief Executive Officer, as the named
sponsor-investigator), formally undertakes the obligations of a sponsor
under 21 CFR Part 312, subpart D, and simultaneously supplies the
cover-sheet metadata (serial number, contents, monitoring and safety
responsibilities) that FDA's document-control system uses to file each
piece of the record. Form 3674 is the instrument through which the
sponsor discharges the separate statutory duty, created by section 402(j)
of the Public Health Service (PHS) Act and enforced through 21 U.S.C.
\S 355 and 42 U.S.C. \S 282(j), to certify the clinicaltrials.gov status
of every clinical investigation referenced in the application. A defect
in either form, an unsigned box 13 on Form 1571 or an unselected
certification box on Form 3674, is a clerical deficiency that can stall
acceptance of an otherwise complete scientific package, which is why this
section treats both forms with the same rigor applied to the protocol and
the Investigator's Brochure \cite{phase1guidance, paisitedocpk}.

\needspace{8\baselineskip}
Figure 3 traces the cover-sheet flow from the Cover Letter through Form
1571 (its serial logic and dual purpose) into Form 3674 (its three
mutually exclusive certification boxes) and onward to the assembled IND
that starts the 30-day FDA review clock under 21 CFR \S 312.40.

\begin{mermaidfig}
\node[mmaccent] (cl)   at (-0.4,0)      {Cover Letter\\sponsor-investigator,\\drug, serial 0000};
\node[mmgoal]   (f1)   at (-0.4,-2.4)   {FDA Form 1571\\(1) sponsor agreement\\(2) cover sheet for\\every submission};
\node[mmin]     (ser)  at (-4.8,-4.9) {Box 10 serial number\\initial = 0000\\(0001, 0002 ... follow)};
\node[mmproc]   (f3)   at (4.4,-4.9)  {FDA Form 3674\\clinicaltrials.gov\\certification of compliance};
\node[mmctx]    (ba)   at (-0.4,-8.2)  {Box A\\no clinical\\investigation};
\node[mmdec]    (bb)   at (4.4,-8.1)  {Box B\\402(j) PHS Act\\does NOT apply\\(typical Phase 1)};
\node[mmctx]    (bc)   at (9.2,-8.2)  {Box C\\402(j) applies;\\will comply};
\node[mmgoal]   (sub)  at (-0.4,-10.8) {Assembled IND\\submitted to FDA;\\30-day review\\clock starts};
\draw[mmedge] (cl) -- (f1);
\draw[mmedge] (f1) to[out=200,in=70,looseness=0.3] (ser.north);
\draw[mmedge] (f1) to[out=340,in=110,looseness=0.3] (f3.north);
\draw[mmedge] (f3.south) to[out=250,in=70,looseness=0.3] (ba.north);
\draw[mmedge] (f3.south) -- (bb.north);
\draw[mmedge] (f3.south) to[out=300,in=110,looseness=0.3] (bc.north);
\draw[mmedge] (bb) to[out=250,in=20,looseness=0.3] node[mmlabel]{selected for this trial} (sub.east);
\draw[mmedge] (ser.south) to[out=-90,in=180,looseness=0.3] (sub.west);
\end{mermaidfig}
\vspace{-0.7cm}
\figcaption{Figure 3. The Cover Letter, FDA Form 1571, and FDA Form 3674
cover-sheet flow. Form 1571 serves the dual purpose of a
sponsor-investigator agreement and a cover sheet for every submission; the
initial serial number in box 10 is 0000. Form 3674 certifies
clinicaltrials.gov compliance through one of Box A (no clinical
investigation), Box B (an investigation to which 42 U.S.C. \S 282(j) /
section 402(j) of the PHS Act does not apply, the typical Phase 1
selection), or Box C (the requirements apply). The Cover Letter and the
FDA forms precede the Table of Contents.}

\subsection{FDA Form 1571 - Investigational New Drug Application}

FDA Form 1571 serves a dual purpose that defines its central place in the
submission. First, it is the binding sponsor-investigator agreement: by
signing box 13, the sponsor-investigator, ChemicalQDevice (Kevin Kawchak,
Chief Executive Officer, as the named sponsor-investigator), commits that
the investigation will not begin until 30 days after FDA receives the IND
(unless notified sooner that the studies may proceed), that a reviewing
Institutional Review Board (IRB) compliant with 21 CFR Part 56 will be
responsible for the initial and continuing review of each study, that the
investigation will be conducted in accordance with all applicable
regulatory requirements, and that an investigator will obtain informed
consent under 21 CFR Part 50 \cite{cfr050paiuni, ichpaistduni}. Second,
and on a recurring basis, Form 1571 is the cover sheet that accompanies
every subsequent submission to the IND file, including protocol
amendments, information amendments, IND safety reports, annual reports,
and general correspondence; each such submission carries its own Form 1571
bearing the next sequential serial number so that the agency can
reconstruct the complete history of the file \cite{phase1guidance}.

The form opens with sponsor identification and contact fields. Box 1
(Name of Sponsor) and boxes 2 through 5 (Address, including street, city,
state, and ZIP) carry the sponsor-investigator entity ChemicalQDevice and
its San Diego, California place of business; box 6 (Telephone Number)
and the associated electronic-contact information complete the routing so
that the assigned project manager and the safety contact can be reached.
Box 7 captures the name(s) of any contract research organization, which
for this single-academic-site study is limited to the named clinical site
and the entities providing investigational product and device support.
The IND Number field is left blank on the initial submission because the
number is assigned by FDA upon receipt; once assigned, that pre-assigned
or receipt-assigned IND number is carried on the Form 1571 of every later
serial submission. Box 9 (Name of Drug) records daraxonrasib (RMC-6236),
the oral RAS(ON) multi-selective small-molecule inhibitor that binds the
GTP-bound active state of mutant RAS in complex with cyclophilin A and
suppresses downstream effector signaling, administered once daily
\cite{daraxwhipple, pdac060s2030}.

The general-information block describes the proposed research and the
phase of investigation. The form indicates that this is a Phase 1
investigation (the first-in-human checkbox), names the indication as
histologically or cytologically confirmed resectable or
borderline-resectable pancreatic ductal adenocarcinoma (PDAC) with a
documented KRAS G12 mutation (G12D, G12V, or G12R) in patients who are
candidates for pancreaticoduodenectomy, and lists the protocol title and
the combined IND and IDE nature of the application. The contents-of-the-
application checkboxes (box 12) are marked to indicate which items of 21
CFR \S 312.23 are included in this serial submission: the Table of
Contents, the Introductory Statement and General Investigational Plan, the
Investigator's Brochure, the study protocol(s), Chemistry, Manufacturing,
and Controls (CMC) information, pharmacology and toxicology information,
previous human experience, and additional information. Box 11 designates
the individuals responsible for monitoring the conduct and progress of the
clinical investigation and the person(s) responsible for review and
evaluation of information relevant to the safety of the drug; for this
study those responsibilities map onto the three-tier oversight structure
(the Data and Safety Monitoring Board, the Independent Safety Monitor, and
the Physical AI Safety Review Committee) described in the protocol
\cite{paiautospons, clintrustpai}.

Because the present application is a combined drug-device study, the box 12
contents block carries an additional editorial burden relative to a
drug-only Phase 1 IND. The CMC item must cross-reference not only the drug
substance and drug product modules for daraxonrasib but also the
device-description and engineering-control material that establishes the
on-premises eight-arm robotic Whipple platform as a significant-risk
investigational device, including its 56 degrees of freedom (eight arms at
seven degrees of freedom each), its 640 sensor channels (80 per arm), its
positional accuracy of 0.05 mm RMS at 1200 mm/s, and its layered force
limits (no more than 3 N per arm and no more than 18 N cumulative). The
pharmacology and toxicology item must reference the computational evidence
base (the 60-second PDAC Whipple-plus-daraxonrasib simulation, the
artificial-intelligence digital-twin PDAC simulation, the quantitative
systems pharmacology metastatic-PDAC simulation with verification,
validation, and uncertainty quantification, and the 100,000-patient in
silico Phase 3 five-arm PDAC study) that supports the integrated
risk assessment \cite{fdadigtwinpc, qspmetpancre, chatgpt100kp,
pdacdigtwinp}. Marking box 12 accurately therefore tells the reviewer, in
a single glance, that the file is internally complete across both the drug
and the device dimensions of the study.

Box 10 is the serial number. FDA requires every submission to an IND
to be numbered consecutively, and the initial submission is serial number
0000. Each subsequent submission to this IND file is numbered
sequentially as 0001, 0002, 0003, and so forth, in the chronological order
in which the agency receives it, regardless of the submission type. For
this initial IND the serial number is therefore 0000, and that value
appears both in box 10 of Form 1571 and in the Cover Letter so that a
reviewer can immediately confirm that this is the originating submission
that starts the 30-day review period under 21 CFR \S 312.40
\cite{phase1guidance, cfr312paiuni}. Table 1.1 summarizes the principal
Form 1571 fields and the value or status of each for this initial
submission.

\needspace{14\baselineskip}
\begin{center}
\begin{xltabular}{\textwidth}{>{\raggedright\arraybackslash}p{4.0cm} Y >{\raggedright\arraybackslash}p{5.5cm}}
\toprule
\textbf{Form 1571 field} & \textbf{Description and content} & \textbf{Value / status (serial 0000)} \\
\midrule
\endfirsthead
\toprule
\textbf{Form 1571 field} & \textbf{Description and content} & \textbf{Value / status (serial 0000)} \\
\midrule
\endhead
\bottomrule
\endfoot
Boxes 1 to 6: Sponsor and contact & Sponsor-investigator name, mailing address, and telephone for the responsible entity. & ChemicalQDevice; San Diego, California; sponsor-investigator Kevin Kawchak, CEO. \\
IND Number & Assigned by FDA on receipt; carried on every later serial submission. & Blank on initial submission (to be assigned). \\
Box 7: CRO / site & Contract research organizations and clinical site(s) supporting the study. & Single academic clinical site; no separate CRO. \\
Box 9: Name of Drug & Investigational drug and active ingredient. & Daraxonrasib (RMC-6236), oral RAS(ON) pan-RAS inhibitor. \\
Phase / indication & Phase of investigation and the disease studied. & Phase 1, first-in-human; resectable or borderline resectable KRAS G12 PDAC. \\
Box 10: Serial Number & Consecutive numbering of every submission to the file. & 0000 (initial); 0001, 0002 ... follow. \\
Box 11: Monitoring / safety & Persons responsible for monitoring conduct and for safety review. & DSMB, Independent Safety Monitor, Physical AI Safety Review Committee. \\
Box 12: Contents & Checkboxes for the 21 CFR \S 312.23 items included. & Introductory statement, general plan, IB, protocol, CMC, pharm/tox, prior human experience, additional information. \\
Box 13: Sponsor agreement & Signed commitments (30-day wait, IRB review, GCP, informed consent). & Signed by sponsor-investigator; executed PDF assembled with submission. \\
\bottomrule
\end{xltabular}
\end{center}
\vspace{-1.75cm}
\figcaption{Table 1.1. Field-by-field summary of FDA Form 1571 for the
initial IND (serial number 0000). The executed\\and signed PDF of the form
is assembled with the physical submission and is not reproduced in this
source.}

Box 13 deserves separate treatment because it is the legally operative
clause of the entire form. The four commitments enumerated there map
directly onto specific provisions of 21 CFR Part 312 and Part 50, and the
sponsor-investigator's signature converts them from regulatory text into a
personal undertaking enforceable by the agency. The 30-day wait
corresponds to 21 CFR \S 312.40(b), which prohibits initiation of the
clinical investigation until 30 days after FDA receives the IND unless the
agency notifies the sponsor that studies may begin earlier; for this
first-in-human Physical AI study the sponsor does not seek early
initiation and will hold all patient-facing activity until the full 30-day
period has elapsed and the Phase 0 simulation-validation gate (a unified
safety level of at least 7.0, at least 1000 simulations across at least
two frameworks, sim-to-real trajectory under 2 mm and tip-force under
0.5 N) has been satisfied \cite{cfr312paiuni, paipredict4x}. The IRB
commitment corresponds to 21 CFR Part 56 and is reinforced by the
protocol's reliance on a single reviewing IRB. The general-compliance
commitment incorporates the good clinical practice expectations of ICH
E6(R3), and the informed-consent commitment corresponds to 21 CFR Part 50
as supplemented by the Physical AI opt-out provision the protocol grounds
in 21 CFR \S 312.60(f) \cite{cfr050paiuni, ichpaistduni}. Table 1.2 sets
out the four box-13 commitments alongside their codified anchors and the
study-specific manner of discharge.

\needspace{12\baselineskip}
\begin{center}
\begin{xltabular}{\textwidth}{>{\raggedright\arraybackslash}p{4.0cm} >{\raggedright\arraybackslash}p{3.0cm} Y}
\toprule
\textbf{Box 13 commitment} & \textbf{Codified anchor} & \textbf{How discharged in this study} \\
\midrule
\endfirsthead
\toprule
\textbf{Box 13 commitment} & \textbf{Codified anchor} & \textbf{How discharged in this study} \\
\midrule
\endhead
\bottomrule
\endfoot
Investigation will not begin before 30 days after FDA receipt. & 21 CFR \S 312.40(b) & No early-initiation request; all patient-facing activity held until the 30-day clock elapses and the Phase 0 simulation gate passes. \\
A Part 56 IRB will perform initial and continuing review. & 21 CFR Part 56 & Single reviewing IRB engaged for the academic site; continuing review at the protocol-specified cadence with cohort-boundary checkpoints. \\
The investigation will follow all applicable requirements. & 21 CFR Part 312; ICH E6(R3) & Three-tier oversight (DSMB, Independent Safety Monitor, Physical AI Safety Review Committee), binding halt rules enforce ongoing compliance. \\
An investigator will obtain informed consent. & 21 CFR Part 50; \S 312.60(f) & Consent under Part 50 and ICH E6(R3) with a Physical AI opt-out, therapeutic-misconception language, and staged-autonomy disclosure. \\
\bottomrule
\end{xltabular}
\end{center}
\vspace{-1.75cm}
\figcaption{Table 1.2. The four box 13 sponsor-investigator commitments,
their codified anchors in 21 CFR, and\\the manner in which each is
discharged in the LLM-directed PDAC robotic daraxonrasib Phase 1 study.}

The dual purpose of Form 1571 has a practical consequence for the
lifecycle of this Physical AI study. Because each protocol amendment,
information amendment, and IND safety report under 21 CFR \S 312.32 will
itself be filed under a new serial number on its own Form 1571, the form
becomes the running index of the entire investigational record, including
the Physical AI specific reporting events. The protocol defines six
Physical AI reporting triggers under 21 CFR \S 312.32(g) (a serious
Physical AI adverse event reported on the 7-day or 15-day clock; a
system-drug interaction within 15 days; a cybersecurity incident within 15
days, or within 7 days where there is potential for patient harm;
sustained AI model degradation over at least three consecutive procedures
or 24 cumulative hours; sim-to-real divergence exceeding twice tolerance,
that is a trajectory deviation greater than 4 mm or a force deviation
greater than 1.0 N; and a digital-twin discrepancy), and each such report
will enter the file as a serially numbered submission cover-sheeted by
Form 1571 \cite{paiautospons, congress9510}. In this way the form binds
the administrative numbering scheme to the safety-surveillance scheme of
the trial.

A worked example makes the serial-numbering discipline concrete. Suppose
the initial IND (serial 0000) clears the 30-day review and the study
opens; the first protocol amendment that adds a clarifying eligibility
criterion would be filed as serial 0001, an information amendment updating
the Investigator's Brochure as serial 0002, and the first IND safety
report (for instance, a fatal or life-threatening suspected adverse
reaction filed on the 7 calendar day clock under 21 CFR \S 312.32(c)(1)(i))
as serial 0003, each cover-sheeted by its own Form 1571 carrying the
FDA-assigned IND number and the next consecutive serial. If a sustained
AI model degradation trigger fires (model performance below threshold over
at least three consecutive procedures or 24 cumulative hours), the
resulting Physical AI safety report under \S 312.32(g) would take the next
serial in turn. The annual report required by 21 CFR \S 312.33 would
likewise be a serially numbered Form 1571 submission. Because the serial
sequence is strictly chronological and type-agnostic, a reviewer can
reconstruct the full procedural history of the file, drug and device
events interleaved, simply by reading the serial numbers in order
\cite{phase1guidance, paisitedocpk}.

The recurring cover-sheet role of Form 1571 should not be confused with
FDA Form 1572 (Statement of Investigator), which is a separate document
addressed elsewhere in this IND under proposed clinical research. Form
1571 binds the sponsor and indexes the file; Form 1572 binds each clinical
investigator to the protocol and to the obligations of 21 CFR \S 312.60.
For this single-site, sponsor-investigator study the two roles converge in
practice on the same organization, but the forms remain distinct
instruments with distinct signatories and distinct regulatory functions,
and conflating them is a common cause of refuse-to-file findings that this
submission is structured to avoid \cite{phase1guidance, paiautospons}.

\subsection{FDA Form 3674 - Certification of Compliance}

FDA Form 3674 is the certification, required by section 402(j) of the
Public Health Service (PHS) Act (codified at 42 U.S.C. \S 282(j)), that
the sponsor has complied with the registration and results-reporting
requirements of clinicaltrials.gov to the extent those requirements apply
to the submission. The form presents three mutually exclusive
certification boxes, exactly one of which must be selected, and the
selection determines what the sponsor is affirming about the
clinicaltrials.gov status of the studies referenced in the application
\cite{phase1guidance}. The logic of the three boxes is summarized in
Table 1.3.

\needspace{12\baselineskip}
\begin{center}
\begin{xltabular}{\textwidth}{>{\raggedright\arraybackslash}p{1.2cm} Y >{\raggedright\arraybackslash}p{5.5cm}}
\toprule
\textbf{Box} & \textbf{What the box certifies} & \textbf{Applicability to this Phase 1} \\
\midrule
\endfirsthead
\toprule
\textbf{Box} & \textbf{What the box certifies} & \textbf{Applicability to this Phase 1} \\
\midrule
\endhead
\bottomrule
\endfoot
Box A & The submission does not reference any clinical trial; that is, there is no clinical investigation at all. & Not selected. This submission references a Phase 1 clinical investigation. \\
Box B & The submission references a clinical trial, but the requirements of 42 U.S.C. \S 282(j) (section 402(j) of the PHS Act) do not apply to that trial. & Selected. This is a Phase 1 investigation; a Phase 1 trial of a drug is excluded from the "applicable clinical trial" definition that triggers mandatory 402(j) registration and results reporting. \\
Box C & The submission references a clinical trial that is an applicable clinical trial to which 42 U.S.C. \S 282(j) does apply, and sponsor certifies it has complied (or will comply) with registration, results reporting. & Not selected for this initial Phase 1. Would apply to a later applicable (for example, Phase 2 or later controlled) trial. \\
\bottomrule
\end{xltabular}
\end{center}
\vspace{-1.75cm}
\figcaption{Table 1.3. The Box A / B / C certification logic of FDA Form
3674 and the selection for this initial\\Phase 1 IND. Exactly one box is
checked; Box B is selected with the rationale stated in the text.}

For this initial Phase 1 IND, Box B is selected. The rationale is that
the study referenced in this application is a Phase 1, first-in-human
investigation of daraxonrasib together with the LLM-directed robotic
Whipple device, and a Phase 1 trial of a drug falls outside the
statutory definition of an "applicable clinical trial" that triggers the
mandatory clinicaltrials.gov registration and results-reporting
obligations of 42 U.S.C. \S 282(j). Under that definition, the
controlled clinical investigations of a drug or biologic that are subject
to mandatory registration are generally those other than Phase 1 trials;
because the present study is Phase 1 (up to n = 18 treated subjects, with
approximately 36 screened, at a single academic site, using a 3+3 dose
escalation across dose levels of 160 mg, 220 mg, and 300 mg once daily
with a 28-day dose-limiting-toxicity window), the section 402(j)
requirements do not apply to it, and Box B is therefore the correct and
typical Phase 1 selection \cite{phase1guidance, oncotrialllm}.

It is worth stating precisely why the "applicable clinical trial" test
resolves to Box B here, because the device dimension of this study could
otherwise invite confusion. The section 402(j) definition reaches two
families of investigations: applicable drug clinical trials, which are
controlled clinical investigations (other than Phase 1 investigations) of
a drug subject to section 505 of the Federal Food, Drug, and Cosmetic Act;
and applicable device clinical trials, which are prospective clinical
studies of health outcomes comparing an intervention against a control in
a device subject to section 510(k), 515, or 522, together with certain
pediatric postmarket surveillance. The present study is a single-arm,
open-label, non-comparative early-feasibility study with no control arm
and no comparison of health outcomes between groups; it does not meet the
"controlled clinical investigation comparing an intervention with a
control" element of either family, and its drug component is expressly a
Phase 1 investigation excluded from the drug-trial definition. On both the
drug axis and the device axis, therefore, the trial falls outside the
mandatory-registration perimeter, and Box B is the legally accurate
selection rather than a mere convention \cite{phase1guidance,
clintrustpai}. Table 1.4 records this analysis as a compact decision
trail.

\needspace{12\baselineskip}
\begin{center}
\begin{xltabular}{\textwidth}{>{\raggedright\arraybackslash}p{5.5cm} >{\raggedright\arraybackslash}p{2cm} Y}
\toprule
\textbf{Definitional element of 402(j)} & \textbf{Met by this study?} & \textbf{Reason} \\
\midrule
\endfirsthead
\toprule
\textbf{Definitional element of 402(j)} & \textbf{Met by this study?} & \textbf{Reason} \\
\midrule
\endhead
\bottomrule
\endfoot
Controlled investigation of a drug, other than Phase 1. & No & Drug component is expressly a Phase 1, first in human investigation, enumerated exclusion. \\
Controlled device study comparing an intervention with a control on health outcomes. & No & Single-arm, open-label, non-comparative early-feasibility design with no control group and no between-group outcome comparison. \\
Applicable pediatric postmarket device surveillance. & No & Adult population (age $\geq$ to 18); no postmarket surveillance mandate is involved. \\
Net result: applicable clinical trial under 402(j)? & No & Box B (requirements do not apply) is the legally accurate selection. \\
\bottomrule
\end{xltabular}
\end{center}
\vspace{-1.75cm}
\figcaption{Table 1.4. Decision trail showing why the study is not an
"applicable clinical trial" under 42 U.S.C.\\ \S 282(j) on either the drug
axis or the device axis, supporting the Box B selection on Form 3674.}

Selecting Box B does not signify that the trial will go unregistered.
The sponsor records that, notwithstanding the inapplicability of the
mandatory 402(j) requirements, voluntary registration on
clinicaltrials.gov is intended and is encouraged as a matter of
transparency, and that any human clinical trial supported in whole or in
part by National Institutes of Health (NIH) funding is subject to the NIH
policy on clinical-trial registration and results information submission,
which independently requires registration on clinicaltrials.gov
regardless of the Form 3674 box selected. The first-in-human Physical AI
nature of this study, and the public-trust considerations attached to an
LLM-directed surgical device operating under verification-before-
generation safeguards aligned with H.R. 9510 (the Verification Before
Generation in Physical AI Oncology Trials Act of 2026), further support
voluntary, prospective registration so that the design, the staged
autonomy limits, and the safety endpoints are publicly visible before the
first dose \cite{congress9510, clintrustpai, paipredict4x}. 

Registering
the design prospectively also exposes, to public scrutiny, the binding
halt rules (DSMB review and potential halt on a device or procedure
serious adverse event in at least one of three subjects within a cohort,
or on two device-related deaths at any time) and the three co-primary
endpoints (device and procedure serious-adverse-event incidence through 30
days, the maximum tolerated dose and recommended Phase 2 dose of
daraxonrasib via the 3+3 design, and technical feasibility), so that
trial's pre-specified stopping logic cannot be revised silently after
results \cite{paiautospons, oncotrialllm}.

The voluntary-registration posture interacts with the program's planned
trajectory in a way that the sponsor anticipates on the face of this
submission. Should the program advance to a later applicable clinical
trial (for example, a randomized or otherwise controlled Phase 2 study of
daraxonrasib with the robotic platform), that later study would meet the
"controlled clinical investigation other than Phase 1" element of the
402(j) drug-trial definition, Box C would then be selected on that
submission's Form 3674, and the corresponding clinicaltrials.gov
registration (within 21 days of enrolling the first subject) and results
reporting (generally within one year of the primary completion date) would
be certified. Maintaining a single, continuously updated
clinicaltrials.gov record from the Phase 1 stage forward therefore eases
the eventual transition from voluntary Box B registration to mandatory Box
C certification without a gap in the public record \cite{phase1guidance,
paisitedocpk}. Table 1.5 contrasts the present Phase 1 posture with the
anticipated later-phase posture across the dimensions a reviewer is most
likely to check.

\needspace{12\baselineskip}
\begin{center}
\begin{xltabular}{\textwidth}{>{\raggedright\arraybackslash}p{3.6cm} Y Y}
\toprule
\textbf{Dimension} & \textbf{This Phase 1 IND (serial 0000)} & \textbf{Anticipated later applicable trial} \\
\midrule
\endfirsthead
\toprule
\textbf{Dimension} & \textbf{This Phase 1 IND (serial 0000)} & \textbf{Anticipated later applicable trial} \\
\midrule
\endhead
\bottomrule
\endfoot
Form 3674 box & Box B (requirements do not apply). & Box C (requirements apply; compliance certified). \\
Trigger for registration & Voluntary, plus NIH policy if NIH-funded. & Mandatory under 42 U.S.C. \S 282(j). \\
Registration timing & Prospective, before first dose, as a transparency measure. & Within 21 days of first-subject enrollment. \\
Results reporting & Voluntary; supports public trust in the Physical AI design. & Generally within one year of primary completion date. \\
Design feature driving the result & Phase 1, single-arm, non-comparative early feasibility. & Controlled or comparative design (for example, randomized Phase 2). \\
\bottomrule
\end{xltabular}
\end{center}
\vspace{-1.75cm}
\figcaption{Table 1.5. Form 3674 posture for the present Phase 1 IND
contrasted with the anticipated posture for a later applicable clinical
trial, showing the planned transition from a voluntary Box B registration
to a mandatory Box C.}

As with Form 1571, the certification on Form 3674 is provided on the
executed, signed PDF that is assembled with the physical submission; the
signature attests that the information on the certification is true and
complete and acknowledges that a knowing and willful false certification
is a criminal offense under 18 U.S.C. \S 1001, and that a false
certification may also subject the certifier to civil monetary penalties
and to the withholding or recovery of grant funds under the enforcement
provisions of 42 U.S.C. \S 282(j)(5). Because Form 3674 must
accompany an IND submission that references one or more clinical trials,
a current Form 3674 is filed with the initial IND (serial 0000) and is
re-filed, with the box selection re-evaluated, whenever a later
submission adds or materially changes a referenced clinical
investigation, keeping the clinicaltrials.gov certification synchronized
with the evolving scope of the program \cite{phase1guidance,
paisitedocpk}. In practice the sponsor treats any amendment that
introduces a new protocol, converts the single-arm design to a controlled
design, or adds a comparator as an event that requires a fresh Form 3674
and a fresh evaluation of the Box A / B / C selection, mirroring the
serial-numbering discipline that Form 1571 imposes on the file.

Taken together, the executed Form 1571 and Form 3674, placed ahead of the
Table of Contents with the Cover Letter, complete the administrative
envelope that initiates FDA's review of this LLM-directed PDAC robotic
daraxonrasib Phase 1 application. Form 1571 supplies the binding sponsor
agreement and the serial-numbered cover sheet that will index every future
drug and device event in the file; Form 3674 supplies the
clinicaltrials.gov certification, here Box B, with a documented decision
trail and a stated voluntary-registration posture that anticipates the
eventual Box C transition. The fidelity of these two forms, verified
field by field against 21 CFR \S 312.23 and the section 402(j) definition,
is the clerical foundation on which the scientific modules that follow are
allowed to rest \cite{phase1guidance, cfr312paiuni, congress9510}.
